% GENERATED FILE. Edit canonical lessons, then run npm run generate:books.
% canonical-source-hash: fnv1a64:d3b996bff90e0dcf
% canonical-lessons: ZH-C394-chenggong, ZH-C394-shibai, ZH-C394-maoxian, ZH-C394-baifei, ZH-C394-bubian

\chapter{Success or Failure}
\label{ch:zh-success-or-failure}
\hlchaptermodality{394}

\begin{chapteropening}
\textbf{By the end of this chapter:} \emph{I can talk about succeeding, failing, taking a risk, effort going to waste and things staying the same.}

Complete the last lesson of chapter 394: I can talk about succeeding, failing, taking a risk, effort going to waste and things staying the same.
\end{chapteropening}

\section[chénggōng]{chénggōng — to succeed}

\label{lesson:ZH-C394-chenggong}



\begin{quote}
Before the new one: say the Mandarin for at any time at all, then the Mandarin for going smoothly.
\end{quote}

\subsection*{You'll want to know: chénggōng}
\emph{chénggōng} — \textquotedblleft{}to succeed\textquotedblright{}.

\begin{quote}
\emph{rúguǒ nǐ nǔlì jiù huì chénggōng} — If you work hard, you'll succeed.
\end{quote}

Its characters come later; for now it is a word for the ear.

\begin{cousinweb}
\emph{chénggōng} is built from \zh{成} 'to complete' and \emph{gōng} 'merit', both inherited from Middle Chinese, \emph{dzyeng} and \emph{kuwng} (Guangyun, in Baxter's transcription), which glosses the second 'an achievement'.
\end{cousinweb}

\begin{grammarlens}[title={What you've built}]
One more word for saying what would happen.
\end{grammarlens}

\subsection*{Your turn}
\begin{itemize}
\raggedright
  \item \emph{Say it:} \emph{chénggōng}
  \item \emph{Say it:} \emph{chénggōng}, once more
  \item \emph{Say it:} \emph{rúguǒ nǐ nǔlì jiù huì chénggōng} — If you work hard, you'll succeed.
  \item \emph{Recall:} say the Mandarin for at any time at all, then the Mandarin for going smoothly, then say \emph{chénggōng} again
\end{itemize}

\begin{tcolorbox}[breakable,colback=teal!4,colframe=teal!35!black,title={Before you move on}]
What does chénggōng mean? (\textquotedblleft{}To succeed\textquotedblright{}.) Say it once more.
\end{tcolorbox}
\section[shībài]{shībài — to fail}

\label{lesson:ZH-C394-shibai}



\begin{quote}
Before the new one: say the Mandarin for going smoothly, then the Mandarin for to succeed.
\end{quote}

\subsection*{You'll want to know: shībài}
\emph{shībài} — \textquotedblleft{}to fail\textquotedblright{}.

\begin{quote}
\emph{shībài le yě méi guānxi wǒmen zài lái} — It doesn't matter if we fail; we'll try again.
\end{quote}

Its characters come later; for now it is a word for the ear.

\begin{cousinweb}
\emph{shībài} is built from \emph{shī} 'to lose' and \emph{bài} 'to fall apart', both inherited from Middle Chinese, \emph{syit} and \emph{baejH} (Guangyun, in Baxter's transcription), which glosses the second 'to break of itself'.
\end{cousinweb}

\begin{grammarlens}[title={What you've built}]
One more word for saying what would happen.
\end{grammarlens}

\subsection*{Your turn}
\begin{itemize}
\raggedright
  \item \emph{Say it:} \emph{shībài}
  \item \emph{Say it:} \emph{shībài}, once more
  \item \emph{Say it:} \emph{shībài le yě méi guānxi wǒmen zài lái} — It doesn't matter if we fail; we'll try again.
  \item \emph{Recall:} say the Mandarin for going smoothly, then the Mandarin for to succeed, then say \emph{shībài} again
\end{itemize}

\begin{tcolorbox}[breakable,colback=teal!4,colframe=teal!35!black,title={Before you move on}]
What does shībài mean? (\textquotedblleft{}To fail\textquotedblright{}.) Say it once more.
\end{tcolorbox}
\section[màoxiǎn]{màoxiǎn — to take a risk}

\label{lesson:ZH-C394-maoxian}



\begin{quote}
Before the new one: say the Mandarin for to succeed, then the Mandarin for to fail.
\end{quote}

\subsection*{You'll want to know: màoxiǎn}
\emph{màoxiǎn} — \textquotedblleft{}to take a risk\textquotedblright{}.

\begin{quote}
\emph{rúguǒ nǐ bù xiǎng màoxiǎn jiù bié qù} — If you don't want to take the risk, don't go.
\end{quote}

Its characters come later; for now it is a word for the ear.

\begin{cousinweb}
\emph{màoxiǎn} is built from \emph{mào} 'to wade into' and \emph{xiǎn} 'danger', both inherited from Middle Chinese, \emph{mawH} and \emph{xjemX} (Guangyun, in Baxter's transcription), which glosses the first 'to cover' and 'to wade' and the second 'perilous'.
\end{cousinweb}

\begin{grammarlens}[title={What you've built}]
One more word for saying what would happen.
\end{grammarlens}

\subsection*{Your turn}
\begin{itemize}
\raggedright
  \item \emph{Say it:} \emph{màoxiǎn}
  \item \emph{Say it:} \emph{màoxiǎn}, once more
  \item \emph{Say it:} \emph{rúguǒ nǐ bù xiǎng màoxiǎn jiù bié qù} — If you don't want to take the risk, don't go.
  \item \emph{Recall:} say the Mandarin for to succeed, then the Mandarin for to fail, then say \emph{màoxiǎn} again
\end{itemize}

\begin{tcolorbox}[breakable,colback=teal!4,colframe=teal!35!black,title={Before you move on}]
What does màoxiǎn mean? (\textquotedblleft{}To take a risk\textquotedblright{}.) Say it once more.
\end{tcolorbox}
\section[báifèi]{báifèi — to be wasted (of effort)}

\label{lesson:ZH-C394-baifei}



\begin{quote}
Before the new one: say the Mandarin for to fail, then the Mandarin for to take a risk.
\end{quote}

\subsection*{You'll want to know: báifèi}
\emph{báifèi} — \textquotedblleft{}to be wasted (of effort)\textquotedblright{}.

\begin{quote}
\emph{yàoshi nǐ xiànzài bú zuò yǐqián de nǔlì jiù báifèi le} — If you don't do it now, all your work so far will be wasted.
\end{quote}

Its characters come later; for now it is a word for the ear.

\begin{cousinweb}
\emph{báifèi} is built from \zh{白} 'white, blank' and \emph{fèi} 'to use up', both inherited from Middle Chinese, \emph{baek} and \emph{phj+jH} (Guangyun, in Baxter's transcription), which glosses the first 'the colour of the west'. \zh{白} for 'in vain' is not given there.
\end{cousinweb}

\begin{grammarlens}[title={What you've built}]
One more word for saying what would happen.
\end{grammarlens}

\subsection*{Your turn}
\begin{itemize}
\raggedright
  \item \emph{Say it:} \emph{báifèi}
  \item \emph{Say it:} \emph{báifèi}, once more
  \item \emph{Say it:} \emph{yàoshi nǐ xiànzài bú zuò yǐqián de nǔlì jiù báifèi le} — If you don't do it now, all your work so far will be wasted.
  \item \emph{Recall:} say the Mandarin for to fail, then the Mandarin for to take a risk, then say \emph{báifèi} again
\end{itemize}

\begin{tcolorbox}[breakable,colback=teal!4,colframe=teal!35!black,title={Before you move on}]
What does báifèi mean? (\textquotedblleft{}To be wasted (of effort)\textquotedblright{}.) Say it once more.
\end{tcolorbox}
\section[\texorpdfstring{\zh{不变} (búbiàn)}{不变 (búbiàn)}]{\zh{不变} (búbiàn) — unchanged; staying the same}

\label{lesson:ZH-C394-bubian}



\begin{quote}
Before the new one: say the Mandarin for to take a risk, then the Mandarin for to be wasted (of effort).
\end{quote}

\subsection*{You'll want to know: \zh{不变}}
\textbf{\zh{不变}} — \emph{búbiàn} — \textquotedblleft{}unchanged; staying the same\textquotedblright{}.

\begin{quote}
\textbf{\zh{时间不变地点改了}} \emph{shíjiān búbiàn dìdiǎn gǎi le} — The time stays the same; the place has changed.
\end{quote}

Tone: \zh{不} is said bú here, rising, because biàn after it is a falling tone.

\begin{cousinweb}
\textbf{\zh{不变}} (\emph{búbiàn}) is built from \zh{不} 'not' and \zh{变} 'to change', both inherited from Middle Chinese, \emph{pjut} and \emph{pjenH} (Guangyun, in Baxter's transcription).
\end{cousinweb}

\begin{grammarlens}[title={What you've built}]
One more word for saying what would happen.
\end{grammarlens}

\subsection*{Your turn}
\begin{itemize}
\raggedright
  \item \emph{Say it:} \emph{búbiàn}
  \item \emph{Say it:} \emph{búbiàn}, once more
  \item \emph{Say it:} \emph{shíjiān búbiàn dìdiǎn gǎi le} — The time stays the same; the place has changed.
  \item \emph{Recall:} say the Mandarin for to take a risk, then the Mandarin for to be wasted (of effort), then say \emph{búbiàn} again
\end{itemize}

\begin{tcolorbox}[breakable,colback=teal!4,colframe=teal!35!black,title={Before you move on}]
What does \zh{不变} mean? (\textquotedblleft{}Unchanged; staying the same\textquotedblright{}.) Say it once more.
\end{tcolorbox}
